\renewcommand{\Anno}{1998}

\renewcommand{\NumeroProblema}{4}

\renewcommand{\Testo}{Determinare la probabilità di ottenere una
somma minore od uguale a 14 lanciando 3 dadi perfettamente
equilibrati.}

\renewcommand{\Risposta}{$49/54$}

\renewcommand{\Aiutino}{Conviene calcolare la probabilità di
ottenere una somma maggiore od uguale a 15.  Quanti sono i caso
possibili?  Quanti (e quali) sono i casi favorevoli? }

\renewcommand{\SoluzioneUfficiale}{
}

\renewcommand{\SoluzioniAlternative}{
}

\renewcommand{\Osservazioni}{
}

\renewcommand{\Commenti}{
}
% Aggiunta redazionale: la soluzione preesistente non viene modificata.
\providecommand{\SoluzioneRedazionale}{}
\renewcommand{\SoluzioneRedazionale}{
% SSSUP-BOZZA-AI: SoluzioneRedazionale; autore=ChatGPT; data=2026-11-08; stato=DA-VERIFICARE
% SSSUP-PROVENIENZA: AI: elaborazione nuova dal testo
I $6^3=216$ esiti ordinati sono equiprobabili. Conviene contare quelli per cui la somma supera $14$. Se i risultati sono $d_1,d_2,d_3$, poniamo $z_i=6-d_i$. La condizione $d_1+d_2+d_3\geq15$ diventa
\[
z_1+z_2+z_3\leq3,\qquad z_i\geq0.
\]
Il numero di terne di somma $j$ è $\binom{j+2}{2}$; per $j=0,1,2,3$ si ottengono rispettivamente $1,3,6,10$ terne. Tutte rispettano automaticamente $z_i\leq5$. Gli esiti sfavorevoli sono dunque $20$, e la probabilità richiesta è
\[
1-\frac{20}{216}=\frac{49}{54}.
\]
}
